% Supplementary Information of the ERL Letter (D-107): copy of methods_surge.tex with the
% outcome of the rule (D-077) added.
% Methods: wind-setup model and surge screening (A5). Design: D-059, D-064 to D-066;
% PREREG section 3. Every number comes from \metric{...}.

\subsection*{Wind-setup screening}
The altimetry series gives one level for the whole sea. On this flat shelf, wind setup and the inverse-barometer effect move the local level, and with it the waterline, away from that level within hours. We therefore modelled the local departure from the basin level and excluded, sector by sector, scenes whose predicted departure exceeded a threshold chosen on training data.

\emph{Observed departures.} We used every Surface Water and Ocean Topography (SWOT) Level-2 high-rate raster granule at 250~m over the domain, in its latest processing version (version~D, or C where D did not exist)\citep{swot_raster_d,swot_raster_c,swot_pdd_raster}. Pixels were kept if their height quality was good or suspect, their water fraction was at least \metric{params.swot_min_water_frac}, no ice flag was set, and they lay in the open sea: \metric{params.lowstand_label} low-stand water whose surroundings within \metric{params.sea_radius_km}~km are at least \metric{params.sea_min_water_pct}\% water, on a \metric{params.sea_mask_res_km}~km grid. Over desert and salt flats, pixels flagged as suspect gave heights far below the sea. For each granule and \metric{params.block_size_km}~km block (Spatial blocking), the median elevation was kept if it rested on at least \metric{params.setup_min_pass_px} pixels with a robust scatter (normalised median absolute deviation) of at most \metric{params.setup_max_block_nmad_m}~m; granules less than \metric{params.setup_pass_gap_min}~min apart form one pass. The departure is this elevation minus the basin level at the pass time, interpolated as for the scenes. Block-passes more than \metric{params.setup_max_robust_z} robust standard deviations from their block's median were dropped, and a block needed at least \metric{params.setup_min_block_passes} passes. Removing each block's mean departure takes out any constant difference between the SWOT and DAHITI references within the block, together with its mean setup. The de-meaned departures were averaged per sector and pass, separately for training and calibration blocks: \metric{surge.n_passes} passes and \metric{surge.n_block_passes} block-passes in all.

\emph{Model.} Hourly 10~m wind and mean sea-level pressure from ERA5\citep{hersbach2020} were averaged over each sector, land included. The features are the means of $|U|u$ and $|U|v$, a quadratic proxy for the wind stress, over the windows \metric{params.setup_windows_h}~h before the acquisition, and the mean pressure anomaly over the preceding \metric{params.setup_pressure_window_h}~h, standardised on the training data. For each sector, the training departures were regressed on them with a ridge penalty chosen from \metric{params.setup_ridge_n} values between zero and \metric{params.setup_ridge_max} by \metric{params.setup_cv_folds}-fold cross-validation over training passes grouped by calendar month (seed \metric{params.seed}). On calibration blocks, the $R^2$ of the \metric{surge.n_sectors_fitted} fitted sectors ranged from \metric{surge.r2_calib_min} to \metric{surge.r2_calib_max} and the residual standard deviation from \metric{surge.resid_sd_calib_min_m} to \metric{surge.resid_sd_calib_max_m}~m (Supplementary Fig.~\ref{fig:setup}); test blocks were never used. Because the sector means include land, the weaker fits may partly reflect the forcing rather than the sea. Komsomolets Bay, shallow and largely dry at the low stand, has no open-sea SWOT pixel. It uses the model of the adjacent northeastern coast with its own forcing, and its waterline statistics are also reported without screening.

\emph{Threshold.} The threshold $\tau$ on the absolute predicted departure in a scene's own sector was chosen among no screening and \metric{params.tau_candidates_m}~m by the pre-registered rule: the smallest $\tau$ beyond which stricter screening no longer reduces the within-bin scatter of waterlines by more than its bootstrap standard error. The scatter $S(\tau)$ is the median, over the epochs, bins and blocks with waterlines in both acquisition halves, of the median shoreline displacement between the halves (the noise-floor measure of the hindcast audit), using only scenes within $\tau$. It was computed with Landsat scenes on a stratified random sample of training blocks, with standard errors from \metric{params.bootstrap_n} resamples of those blocks. Starting from no screening, $\tau$ was tightened until the next stricter step reduced $S$ by less than that step's standard error, or reached the strictest candidate. No lidar, calibration-block or test-block value enters the choice, which applies to both sensor families. A scene above $\tau$ in a sector is excluded in that sector only; scenes newer than the last ERA5 hour lack forcing and are excluded everywhere. The number of scenes that each candidate removes is reported.

\emph{Outcome.} The rule selected \metric{screening.chosen_tau}. On \metric{screening.n_blocks} sampled training blocks, $S$ was \metric{screening.S_m.inf}~m without screening and stayed at one \metric{params.stage2_res_m}~m grid cell at every threshold (\metric{screening.S_m.0p05}~m at the strictest), while the number of usable block-bins fell from \metric{screening.n_block_bins.inf} to \metric{screening.n_block_bins.0p05} (Supplementary figure~\ref{fig:method}c). The statistic therefore sits at the resolution of the grid, which limits what the rule can detect; wind setup enters the robustness checks instead, as a stricter fixed threshold, a setup correction in place of screening, and screening in Komsomolets Bay (section~\SIrobust).
